\section{Additional benchmark and agent results}
\label{sec:new-results}
This appendix tabulates the DeepSeek V4 Pro column of
Figure~\ref{fig:swectx-crossmodel} (Table~\ref{tab:new-swectx}) and gives the
LoLBench comparison from
Section~\ref{sec:backbones}.
The three-benchmark comparison is tabulated in
Table~\ref{tab:main-results-overview}; additional comparisons appear in
Appendix~\ref{sec:historical}.

\begin{table}[htbp]
\centering
\caption{DeepSeek V4 Pro column of Figure~\ref{fig:swectx-crossmodel}
(300-summary initial pool; the ContextGraph and no-memory values are those of
the main comparison). Bold marks the best automatic memory result; the oracle
receives a known related summary.}
\label{tab:new-swectx}
\begin{tabular}{lr}
\toprule
Method & Resolution rate (\%) \\
\midrule
No memory & 21.4 \\
Random summaries & 27.1 \\
Mem0 & 26.8 \\
ContextGraph & \textbf{35.1} \\
\midrule
Oracle summary & 27.8 \\
\bottomrule
\end{tabular}
% Source: figures/scripts/crossmodel_data.json, DeepSeek V4 Pro column.
\end{table}

\paragraph{LoLBench Python subset.}
We select 20 Python tasks from LoLBench \citep{lolbench2026}, which evaluates
long-horizon development from enhancement proposals in large software
systems. The agent is Codex with GPT-5.5. The initial experience graph is
built from other Python tasks, separate from the 20 evaluation tasks.
During evaluation, ContextGraph adds lessons from completed attempts,
including successes and failures, for subsequent retrieval. Baseline
memory updates follow each method's open-source defaults.

\begin{table}[htbp]
\centering
\caption{Success rates on 20 LoLBench Python tasks with Codex and GPT-5.5,
as reported in Section~\ref{sec:backbones}.
pass@$k$ reports the fraction of tasks solved within $k$ attempts.}
\label{tab:new-codex55}
\begin{tabular}{llr}
\toprule
Method & Metric & Success rate (\%) \\
\midrule
No memory & pass@1 & 5.0 \\
FAISS Flat & pass@1 & 5.0 \\
ExpeL & pass@1 & 5.0 \\
ExpeRepair & pass@1 & 5.0 \\
ReasoningBank & pass@1 & 5.0 \\
ACE & pass@1 & 5.0 \\
Supermemory & pass@1 & 5.0 \\
ContextGraph & pass@1 & \textbf{15.0} \\
\midrule
ContextGraph & pass@3 & \textbf{25.0} \\
\bottomrule
\end{tabular}
\end{table}

\clearpage
\section{Full-context agent-directed memory selection}
\label{sec:agentic-retrieval}
At step $i$ of task $q_t$, the agent uses the complete context $C_{t,i}$: the issue and
conversation, inspected code, edits, actions, and observed failures. It can
search the whole experience pool, read source records, follow links between
strategies and errors, and compare lessons across repositories. Its selection
follows Equation~\ref{eq:agent-selection}: searches return ranked
candidates, and the agent decides which records to inspect and retain, so the
number of selected lessons varies by task. Each memory action $m_j$ and its
result extend the history to $h_{j+1}$ for the next action; at
\textsc{stop}, the kept lessons form $\mathcal M_{t,i}$.

The selected memory contains strategies and warnings, their source context,
and the conditions under which they are useful. Before planning, selection
helps choose an investigation strategy. After an error, the updated context
helps revise the diagnosis and find additional guidance. After each completed
task, the graph incorporates new experience for subsequent tasks.
The pool therefore evolves along the issue-ordered task stream; official
benchmark test results are never part of the agent's context or stored memory.
